\section{Conclusion}
\label{sec:conclusion}
In this paper, we introduced HopRAG, a novel RAG system with a logic-aware retrieval mechanism. HopRAG connects related passages through pseudo-queries, which allows identifying truly relevant passages within multi-hop neighborhoods of indirectly relevant ones, significantly enhancing both the precision and recall of retrieval.

Extensive experiments on multi-hop QA benchmarks, i.e. MuSiQue, 2WikiMultiHopQA, and HotpotQA, demonstrate that HopRAG outperforms conventional RAG systems and state-of-the-art baselines. Specifically, HopRAG achieved over 36.25\% higher answer accuracy and 20.97\% improved retrieval F1 score compared to conventional information retrieval approaches. It highlights the effectiveness of integrating logical reasoning into the retrieval module. Moreover, ablation studies provide insights into the sensitivity of hyperparameters and models, revealing trade-offs between retrieval performance and computational costs. 

HopRAG paves the way toward reasoning-driven knowledge retrieval. \textbf{Future work} involves scaling HopRAG to broader domains beyond QA tasks; optimizing indexing and traversal strategies for more complex scenarios with lower computation costs.

\section*{Acknowledgments}
This work is supported by the National Key R\&D Program of China (2024YFA1014003), National Natural Science Foundation of China (92470121, 62402016), CAAI-Ant Group Research Fund, and High-performance Computing Platform of Peking University.

\section*{Limitations}
Despite the benefits of HopRAG, the current evaluation focuses on multi-hop or multi-document QA tasks. In order to mitigate the risk of performance fluctuations when applying HopRAG to other datasets, we should explore its generalization capabilities across a broader range of domains. Besides, more sophisticated query simulation and edge merging strategies may lead to further improvements. Finally, though we were inspired by the theories of six degrees of separation and small-world networks, the degree distribution of our passage graph vertices does not exhibit the power-law characteristic. On the one hand, these theories only serve as the motivation and intuitive analogy; on the other hand, exploring more appropriate degree distribution strategies is an interesting research topic. We leave these research problems for future work.

\section*{Ethics Impact}
In our work, we acknowledge two key ethical considerations. First, we utilized AI assistant to enhance the writing process of our paper and code. We ensure that the AI assistant was used as a tool to improve clarity and conciseness, while the final content and ideas were developed and reviewed by human authors. Second, we employed multiple open source datasets and one open source tool Neo4j Community Edition \citep{10.1145/2384716.2384777} under GPL v3 license in our experiments. We are transparent about their origin and limitations, and we respect data ownership and user privacy. 


